\documentclass[../../../main.tex]{subfiles}

\begin{document}

\section{Vectors}

We have discussed on vectors more in-depth in the part for Linear
Algebra. In this section however, let's focus more on the geometric
intuition than the abstract nature of vectors.

\begin{definition}{}{}
\textbf{Vectors} are quantities with both magnitude and a direction.
\end{definition}

Some example of vectors include force, velocity, and acceleration
because they all have a specific direction and also a magnitude.
When we graph them, we can think of them as a arrow with starting
and ending points.

\begin{figure}[H]
\centering
\tdplotsetmaincoords{70}{110}
\begin{tikzpicture}[
        tdplot_main_coords,
        axis/.style={->,thick},
        scale=0.4
    ]
	\draw[axis] (-8,0,0) -- (8,0,0) node[left] {$x$};
	\draw[axis] (0,-7,0) -- (0,8,0) node[right] {$y$};
    \draw[axis] (0,0,-5) -- (0,0,8) node[left] {$z$};

    \draw[axis] (2,-3,2) -- (1,3,5) node[above right] {$\mathbf{v}$};
    \draw[axis] (2,-3,-1) -- (1,3,2) node[above right] {$\mathbf{w}$};
    \node[circle,fill=black,inner sep=0pt,minimum size=1mm]
        at (2,-3,2) {};
    \node[circle,fill=black,inner sep=0pt,minimum size=1mm]
        at (2,-3,-1) {};
\end{tikzpicture}
\end{figure}

From now, let's represent vectors with $\mathbf{u}, \mathbf{v},
\mathbf{w}$ and scalars, or real numbers, with $a, b, \alpha, \beta$.
From the graph above, we can define another definition.

\begin{definition}{}{}
The \textbf{terminal point} of a vector is the point where the vector
ends. The \textbf{initial point} is the point where the vector begins.
\end{definition}

When we want to add two vectors $\mathbf{v}$ and $\mathbf{w}$,
we would connect the terminal point of $\mathbf{v}$ by the starting
point of $\mathbf{w}$ and obtain a new vector by connecting the
initial point of $\mathbf{v}$ and the terminal point of $\mathbf{w}$
as shown by the figure below.

\begin{figure}[H]
\centering
\tdplotsetmaincoords{70}{110}
\begin{tikzpicture}[
        tdplot_main_coords,
        axis/.style={->,thick},
        scale=0.4
    ]
	\draw[axis] (-8,0,0) -- (8,0,0) node[left] {$x$};
	\draw[axis] (0,-7,0) -- (0,8,0) node[right] {$y$};
    \draw[axis] (0,0,-5) -- (0,0,8) node[left] {$z$};

    \draw[axis] (2,-3,-1) -- (1,-1,2) node[above] {$\mathbf{v}$};
    \draw[axis] (1,-1,2) -- (1,6,4) node[above] {$\mathbf{w}$};
    \draw[axis, Red] (2,-3,-1) -- (1,6,4)
        node[below right] {$\mathbf{v + w}$};
    \node[circle,fill=black,inner sep=0pt,minimum size=1mm]
        at (2,-3,-1) {};
    \node[circle,fill=black,inner sep=0pt,minimum size=1mm]
        at (1,-1,2) {};
\end{tikzpicture}
\end{figure}

Multiplying a vector by scalar $k$ can be though of as lengthening
the vector by a factor of $k$. If $k < 0$, the vector is pointing
the opposite direction.

\begin{figure}[H]
\centering
\tdplotsetmaincoords{70}{110}
\begin{tikzpicture}[
        tdplot_main_coords,
        axis/.style={->,thick},
        scale=0.4
    ]
    \draw[axis] (-8,0,0) -- (8,0,0) node[left] {$x$};
    \draw[axis] (0,-7,0) -- (0,8,0) node[right] {$y$};
    \draw[axis] (0,0,-5) -- (0,0,8) node[left] {$z$};

    \draw[axis] (2,-3,2) -- (1.5,0,3.5) node[above] {$\mathbf{v}$};
    \draw[axis] (-2,3,-2) -- (-1.5,-0,-3.5)
        node[right] {$\mathbf{-v}$};
    \draw[axis] (2,-3,-1) -- (1,3,2) node[below] {$\mathbf{2v}$};
    \node[circle,fill=black,inner sep=0pt,minimum size=1mm]
        at (2,-3,2) {};
    \node[circle,fill=black,inner sep=0pt,minimum size=1mm]
        at (2,-3,-1) {};
\end{tikzpicture}
\end{figure}

Also for the purpose of the notes, let's consider the vectors
with same magnitude and direction but different initial vectors
as the same vectors. The being said, we can write any vectors
in $n$-dimension as $\langle v_1, v_2, v_3, \ldots, v_n \rangle$
where $(v_1, v_2, v_3, \ldots, v_n)$ is the terminal point and
$(0, 0, \linebreak, 0 \ldots, 0)$ is the initial point.
Such numbers $v_1, \ldots, v_n$ are known as the \textbf{components}
of $\mathbf{v}$.

Now that we have defined $-\mathbf{v}$, we can consider vector
subtraction. We can consider vector subtraction as adding two vectors
where one is scaled by a negative factor.

\begin{figure}[H]
\centering
\tdplotsetmaincoords{70}{110}
\begin{tikzpicture}[
        tdplot_main_coords,
        axis/.style={->,thick},
        scale=0.4
    ]
    \draw[axis] (-8,0,0) -- (8,0,0) node[left] {$x$};
    \draw[axis] (0,-7,0) -- (0,8,0) node[right] {$y$};
    \draw[axis] (0,0,-5) -- (0,0,8) node[left] {$z$};

    \draw[axis] (2,2,-2) -- (1,5,3)
        node[right] {\footnotesize $\mathbf{v}$};
    \draw[axis] (2,2,-2) -- (2,-2,2)
        node[left] {\footnotesize $\mathbf{w}$};
    \draw[axis,dotted] (1,5,3) -- (1,1,7);
    \draw[axis,dotted] (2,-2,2) -- (1,1,7);
    \draw[axis,Red] (2,-2,2) -- (1,5,3)
        node[pos=0.7, below] {\footnotesize $\mathbf{v - w}$};
    \draw[axis,Blue] (2,2,-2) -- (1,1,7)
        node[pos=0.7, left] {\footnotesize $\mathbf{v + w}$};
\end{tikzpicture}
\end{figure}

Using components, we can write addition, subtraction and scalar
multiplication as the following for $\mathbf{v} = \langle v_1, v_2,
v_3 \rangle$ and $\mathbf{w} = \langle w_1, w_2, w_3 \rangle$.
\begin{align*}
    \mathbf{v} + \mathbf{w}
    &= \langle v_1 + w_1, v_2 + w_2, v_3 + w_3 \rangle \\
    \mathbf{v} - \mathbf{w}
    &= \langle v_1 - w_1, v_2 - w_2, v_3 - w_3 \rangle \\
    k\mathbf{v}
    &= \langle kv_1, kv_2, kv_3 \rangle
\end{align*}
As you can see from the equations above, we can see that vector
addition is commutative and associative and scalar multiplication is
distributive. I will keep them brief here because we already discussed
about these in-depth on my other notes for linear algebra. Now this
property leads to the natural question of subtracting a vector by
itself.

\begin{definition}{}{}
The \textbf{zero vector} is a vector with magnitude of zero.
\end{definition}

Intuitively, we can think of a zero vector as vector with the same
initial and terminal points. Another way of thinking about a zero
vector is a vector with zero as its components. This will naturally
leads to $\mathbf{v} + \mathbf{0} = \mathbf{v}$ for any vector
$\mathbf{v}$. For the purpose of this notes, let's consider zero
vector as the vector that is parallel to every vectors.

Continuing with the definition of components and the magnitude of a
vector, we can define an important terminology.

\begin{definition}{}{}
The magnitude of a vector $\mathbf{v}$ is known as the \textbf{norm}
of $\mathbf{v}$, denoted as $\| \mathbf{v} \|$.
\end{definition}

Notice that the magnitude, or length, cannot be negative. Therefore,
we can see that the norm of any vector is non-negative. Moreover,
$\| \mathbf{0} \| = 0$. Also note that scalar multiplication is
defined element-wise. Using the distance formula, we can also obtain
the following very important equation for a vector $\mathbf{v}$ and
scalar $k$.
\[
\| k \mathbf{v} \| = |k| \cdot \| \mathbf{v} \|
\]
Let's take a look at a quick example. Consider the vector
$\mathbf{v} = \langle 1, 2, 3 \rangle$. Notice that the following
equation holds.
\[
\| \mathbf{v} \|
= \sqrt{1^2 + 2^2 + 3^2}
= \sqrt{14}
\]
We can also see that the following equation holds with scalar $2$
and $2 \langle 1, 2, 3 \rangle = \langle 2, 4, 6 \rangle$.
\begin{align*}
    \| 2 \mathbf{v} \|
    &= \sqrt{2^2 + 4^2 + 6^2}
    = \sqrt{2^2 (1^2 + 2^2 + 3^2)} \\
    &= 2 \sqrt{1^2 + 2^2 + 3^2} \\
    &= 2 \cdot \| \mathbf{v} \|
\end{align*}
With this equation in mind, let's consider the following question.

\begin{important}
Given a vector $\mathbf{v}$, how do we find a unit vector, or
vector with length $1$, with the same direction as $\mathbf{v}$?
\end{important}

One way of thinking this is rearranging our equation
$\| k \mathbf{v} \| = |k| \cdot \| \mathbf{v} \|$. Consider the
following equations with unit vector $\mathbf{u}$ and
$\| \mathbf{v} \| = k \neq 0$.
\[
\left| \frac{1}{k} \right| \cdot \| \mathbf{v} \|
= \frac{1}{k} \cdot k
= 1
= \| \mathbf{u} \|
\]
This solves it! We can find a unit vector with the same direction
$\mathbf{v}$ by dividing it by its norm. Another variation of this
equation would be the following.
\[
\| \mathbf{v} \| = |k| \cdot \| \mathbf{u} \|
\]
\begin{important}
This implies that any vector $\mathbf{v}$ can be represented as
the product of its norm and a unique unit vector with the same
direction as $\mathbf{v}$. Such process of multiplying a nonzero
vector by the reciprocal of its norm is known as \textbf{normalizing}.
\end{important}

Finally, another standard way of writing vectors is by linear
combination of the standard basis vectors defined as the following.
\[
\mathbf{i} = \langle 1, 0, 0 \rangle, \quad
\mathbf{j} = \langle 0, 1, 0 \rangle, \quad
\mathbf{k} = \langle 0, 0, 1 \rangle
\]
With this, we can write $\mathbf{v} = \langle 1, 2, 3 \rangle$ as
$\mathbf{i} + 2\mathbf{j} + 3\mathbf{k}$.

More detailed and abstract introduction of vectors are included
in the linear algebra part of LibreNotebook, and I highly recommend
checking that out! From here, we will discuss more foundational and
geometric interpretations of vectors that we have not discussed
in my notes for linear algebra.

\subsection{Dot Products and Projections}

As always, let's start with definition.

\begin{definition}{}{}
The \textbf{dot product} of two vectors $\mathbf{u}$ and $\mathbf{v}$
of the same dimension is defined as the sum of all products of
elements in the same index.
\end{definition}

Writing them with mathematical notations, we can write them as the
following.
\[
\langle u_1, u_2, \ldots, u_n \rangle \cdot
    \langle v_1, v_2, \ldots, v_n \rangle
= \sum_{i=1}^n u_i v_i
\]
Here is a quick example.
\[
\langle 1, 2, 3 \rangle \cdot \langle -1, 0, 1 \rangle
= -1 + 0 + 3
= 2
\]
Continuing, below are some properties of dot products. The statements
below are quite self-evident as they mirror the properties of addition
and subtraction.

\begin{theorem}{}{Properties of the Dot Product}
For vectors $\mathbf{u}, \mathbf{v}, \mathbf{w}$ of the same dimension
and an arbitrary scalar $\lambda$, the following equations hold.
\begin{enumerate}
\item $\mathbf{u} \cdot \mathbf{v} = \mathbf{v} \cdot \mathbf{u}$
\item $\lambda (\mathbf{u} \cdot \mathbf{v})
= (\lambda \mathbf{u}) \cdot \mathbf{v}
= \mathbf{u} \cdot (\lambda \mathbf{v})$
\item $\mathbf{u} \cdot (\mathbf{v} + \mathbf{w})
= \mathbf{u} \cdot \mathbf{v} + \mathbf{u} \cdot \mathbf{w}$
\item $\mathbf{v} \cdot \mathbf{v} = \| v \|^2$
\end{enumerate}
\end{theorem}

One of the reasons why the dot product of two vectors is so important
is because it helps us determine the angle between the two vectors.
Consider the following theorem.

\begin{theorem}{}{The Dot Product and the Angle}
Let $\mathbf{u}, \mathbf{v}$ be two nonzero vectors. Let $\theta$ be
the smaller angle between the two vectors when their initial point
coincide. Then, the following equation holds.
\[
\mathbf{u} \cdot \mathbf{v}
= \| u \| \| v \| \cos\theta
\]
\end{theorem}
\begin{proof}
First and foremost, consider the following equation by the law of
cosines.
\[
\cos\theta
= \frac{ \| \mathbf{u} \|^2 + \| \mathbf{v} \|^2 -
    \| \mathbf{v} - \mathbf{u} \|^2
}{ 2 \| \mathbf{u} \| \| \mathbf{v} \| }
\]
By Theorem \ref{theo:Properties of the Dot Product}, the following
equations hold.
\begin{align*}
    \cos\theta
    &= \frac{ \| \mathbf{u} \|^2 + \| \mathbf{v} \|^2 -
        ( \mathbf{v} - \mathbf{u} ) \cdot ( \mathbf{v} - \mathbf{u} )
    }{ 2 \| \mathbf{u} \| \| \mathbf{v} \| } \\
    &= \frac{ \| \mathbf{u} \|^2 + \| \mathbf{v} \|^2 - \left(
        \mathbf{v} \cdot \mathbf{v}  - 2 \mathbf{u} \cdot \mathbf{v}
            + \mathbf{u} \cdot \mathbf{u}
        \right)
    }{ 2 \| \mathbf{u} \| \| \mathbf{v} \| } \\
    &= \frac{2 \mathbf{u} \cdot \mathbf{v}}
        { 2 \| \mathbf{u} \| \| \mathbf{v} \| } \\
    \therefore 
    \| u \| \| v \| \cos\theta
    &= \mathbf{u} \cdot \mathbf{v}
\end{align*}
Thus, the theorem holds.
\end{proof}

Now this theorem tells us the relationship between the angle between
two vectors and their dot product.

\begin{important}
If $0 \leq \theta < \frac{\pi}{2}$, then $\mathbf{u} \cdot \mathbf{v}
> 0$. If $\theta = \frac{\pi}{2}$, then $\mathbf{u} \cdot \mathbf{v}
= 0$. If $\frac{\pi}{2} < \theta < \pi$, then $\mathbf{u} \cdot
\mathbf{v} < 0$.
\end{important}

This somewhat simple observation with the cosine function and the
fact that $\| \mathbf{u} \|, \| \mathbf{u} \| > 0$ reveals the
relationship between the dot product of two vectors and their angles.
Continuing from this observation, we can establish a definition.

\begin{definition}{}{}
Two nonzero vectors $\mathbf{u}$ and $\mathbf{v}$ are
\textbf{orthogonal} if they form a right angle. In other words,
$\mathbf{u} \cdot \mathbf{v} = 0$.
\end{definition}

Continuing with the relationship and definition, let's discuss
more on the angle of vectors.

\begin{definition}{}{}
The \textbf{direction angles} of a nonzero vector $\mathbf{v}$ is
the angle between $\mathbf{v}$ and the standard basis vectors.
\end{definition}

To find the direction angles of $\mathbf{v} = \langle v_1, v_2, v_3
\rangle$, let's return to Theorem \ref{theo:The Dot Product and the
Angle}. Consider the following equation for the angle $\alpha$
between $\mathbf{v}$ and the standard basis vector $\mathbf{i}$.
\begin{align*}
    \mathbf{v} \cdot \mathbf{i}
    &= \| \mathbf{v} \| \| \mathbf{i} \| \cos\alpha \\
    v_1 &= \| \mathbf{v} \| \cos\alpha \\
    \therefore
    \alpha &= \cos^{-1} \left( \frac{v_1}{\| \mathbf{v} \|} \right)
\end{align*}
Following the same pattern for $\beta$ and $\gamma$, the can
establish the following equations.
\[
\cos\alpha = \frac{v_1}{\| \mathbf{v} \|}, \quad
\cos\beta = \frac{v_2}{\| \mathbf{v} \|}, \quad
\cos\gamma = \frac{v_3}{\| \mathbf{v} \|}
\]
Indeed, we have special term for the values above. The values above
are known as the \textbf{direction cosines} of $\mathbf{v}$. For
vectors of direction cosines, the can write the following equation.
\[
\langle \cos\alpha, \cos\beta, \cos\gamma \rangle
= \left\langle
    \frac{v_1}{\| \mathbf{v} \|},
    \frac{v_2}{\| \mathbf{v} \|},
    \frac{v_3}{\| \mathbf{v} \|}
\right\rangle
= \frac{\mathbf{v}}{\| \mathbf{v} \|}
\]
Continuing, let's consider the case when $\| \mathbf{v} \| = 1$.
In that case,
\[
\langle \cos\alpha, \cos\beta, \cos\gamma \rangle
= \left\langle
    \frac{v_1}{\| \mathbf{v} \|},
    \frac{v_2}{\| \mathbf{v} \|},
    \frac{v_3}{\| \mathbf{v} \|}
\right\rangle
= \mathbf{v}
\]
also hold! Finally, here is another interesting properties of the
direction cosines.
\[
\cos^2\alpha + \cos^2\beta + \cos^2\gamma
= \left( \frac{v_1}{\| \mathbf{v} \|} \right)^2
    + \left( \frac{v_2}{\| \mathbf{v} \|} \right)^2
    + \left( \frac{v_3}{\| \mathbf{v} \|} \right)^2
= \frac{\| \mathbf{v} \|^2}{\| \mathbf{v} \|^2}
= 1
\]
We can see that for nonzero $\mathbf{v}$, we can generalize this
for any $n$. Therefore the sum of the squares of direction cosines
are always $1$.

The last part that I wish to discuss about dot products is
projections. First, consider a vector $\mathbf{v} = \mathbf{w}_1
+ \mathbf{w}_2$ in a plane where $\mathbf{w}_1$ and $\mathbf{w}_2$ are
along the standard basis $\mathbf{i}$ and $\mathbf{j}$ respectively.
By construction, we can write $\mathbf{w}_1 = \mathbf{v} \cdot
\mathbf{i}$ and $\mathbf{w}_2 = \mathbf{v} \cdot \mathbf{j}$.

\begin{important}
In other words, for any nonzero vector $\mathbf{v}$ in a plane,
it can be represented as the following.
\[
\mathbf{v}
= (\mathbf{v} \cdot \mathbf{i}) \mathbf{i} +
    (\mathbf{v} \cdot \mathbf{j}) \mathbf{j}
\]
Such vectors $(\mathbf{v} \cdot \mathbf{i}) \mathbf{i}$ and
$(\mathbf{v} \cdot \mathbf{j}) \mathbf{j}$ are known as
\textbf{vector components} of $\mathbf{v}$ along $\mathbf{i}$ and
$\mathbf{j}$ where the dot products $\mathbf{v} \cdot \mathbf{i}$ and
$\mathbf{v} \cdot \mathbf{j}$ are \textbf{scalar components} of
$\mathbf{v}$ along $\mathbf{i}$ and $\mathbf{j}$.
\end{important}

This idea can be generalized for higher dimensions. Moreover,
although we used standard basis vectors here, this works for any
orthogonal unit vectors.

Another way of expanding this is through geometric interpretation.

\begin{minipage}{0.48\textwidth}
\begin{figure}[H]
\centering
\begin{tikzpicture}
    \coordinate (O) at (0,0);
    \coordinate (A) at (4,0);
    \coordinate (B) at (0,3);
    \coordinate (C) at (4,3);
    \coordinate (e1) at (1,0);
    \coordinate (e2) at (0,1);

    \draw[->, thick] (O) -- (A);
    \draw[->, thick] (O) -- (B);
    \draw[->, thick, BrickRed] (O) -- (C);
    \draw[->, thick, RoyalPurple] (O) -- (e1);
    \draw[->, thick, RoyalPurple] (O) -- (e2);
    \draw[dashed, thick] (A) -- (C);
    \draw[dashed, thick] (B) -- (C);

    \pic[draw,"$\theta$",angle radius=8mm,angle eccentricity=1.4]
        {angle=A--O--C};
    \node[below] at (e1) {$\mathbf{e}_1$};
    \node[left] at (e2) {$\mathbf{e}_2$};
    \node[below] at (A) {$\mathbf{w}_1$};
    \node[left] at (B) {$\mathbf{w}_2$};
    \node[above right] at (C) {$\mathbf{v}$};
    \node[below] at ($(B)!0.3!(C)$)
        {\footnotesize $\| \mathbf{v} \| \cos\theta$};
    \node[left] at ($(A)!0.3!(C)$)
        {\footnotesize $\| \mathbf{v} \| \sin\theta$};
\end{tikzpicture}
\end{figure}
\end{minipage}
\hfill
\begin{minipage}{0.48\textwidth}
\begin{figure}[H]
\centering
\begin{tikzpicture}
    \coordinate (O) at (0,0);
    \coordinate (A) at (-4,0);
    \coordinate (B) at (0,3);
    \coordinate (C) at (-4,3);
    \coordinate (e1) at (1,0);
    \coordinate (e2) at (0,1);

    \draw[->, thick] (O) -- (A);
    \draw[->, thick] (O) -- (B);
    \draw[->, thick, BrickRed] (O) -- (C);
    \draw[->, thick, RoyalPurple] (O) -- (e1);
    \draw[->, thick, RoyalPurple] (O) -- (e2);
    \draw[dashed, thick] (A) -- (C);
    \draw[dashed, thick] (B) -- (C);

    \pic[draw,"$\theta$",angle radius=4mm,angle eccentricity=1.6]
        {angle=e1--O--C};
    \node[below] at (e1) {$\mathbf{e}_1$};
    \node[left] at (e2) {$\mathbf{e}_2$};
    \node[below] at (A) {$\mathbf{w}_1$};
    \node[right] at (B) {$\mathbf{w}_2$};
    \node[above] at (C) {$\mathbf{v}$};
    \node[below] at ($(B)!0.4!(C)$)
        {\footnotesize $\| \mathbf{v} \| \cos\theta$};
    \node[right] at ($(A)!0.3!(C)$)
        {\footnotesize $\| \mathbf{v} \| \sin\theta$};
\end{tikzpicture}
\end{figure}
\end{minipage}

As you can see from the diagrams above, we can write the scalar
components with respect to the norm of $\mathbf{v}$ and the
angle $\theta$. Moreover, we can define a new definition.

\begin{definition}{}{}
The \textbf{orthogonal projection} of a vector $\mathbf{v}$,
denoted as $\proj_\mathbf{u} \mathbf{v}$, on a unit vector
$\mathbf{u}$ is the vector component of $\mathbf{v}$ that is parallel
to $\mathbf{u}$.
\end{definition}

From the figures above, we can see that $\proj_{\mathbf{e}_1} =
\mathbf{w}_1$ and $\proj_{\mathbf{e}_2} = \mathbf{w}_2$.
Finally, this leads us to the conclusion that for any vector
$\mathbf{v}$ and respective orthogonal unit vectors $\mathbf{e}_1$
and $\mathbf{e}_2$, the following equation holds.
\[
\mathbf{v}
= \proj_{\mathbf{e}_1} \mathbf{v} + \proj_{\mathbf{e}_2} \mathbf{v}
\]
This is it for dot products and orthogonal projections! Before
we conclude our section on vectors, we will discuss about cross
products and its geometric interpretations.

\subsection{The Cross Product and Geometry}

Unlike how we defined vector addition and scalar multiplication,
we had a special way of defining ``multiplication'' of vectors
with the dot product. However specifically for $3$-dimensional
space, we can also define cross products.

\begin{definition}{}{}
For $\mathbf{u} = \langle u_1, u_2, u_3 \rangle$ and $\mathbf{v} =
\langle v_1, v_2, v_3 \rangle$, the \textbf{cross product},
denoted as $\mathbf{u} \times \mathbf{v}$ and read as
``$\mathbf{u}$ cross $\mathbf{v}$'', is defined as the following.
\vspace{-0.8\baselineskip}
\[
\mathbf{u} \times \mathbf{v}
= (u_2v_3 - u_3v_2) \mathbf{i} - (u_1v_3 - u_3v_1) \mathbf{j}
    + (u_1v_2 - u_2v_1) \mathbf{k}
\]
\end{definition}

As you can see from the definition above, unlike the dot product
that returns a scalar, the cross product of two vector return
a vector. Also, if it is difficult to memorize the formula above,
we can also think of cross product as the following.
\begin{align*}
    \mathbf{u} \times \mathbf{v}
    &= \det\left(
        \begin{bmatrix}
            \mathbf{i} & \mathbf{j} & \mathbf{k} \\
            u_1        & u_2        & u_3        \\
            v_1        & v_2        & v_3
        \end{bmatrix}
    \right) \\
    &= (u_2v_3 - u_3v_2) \mathbf{i} - (u_1v_3 - u_3v_1) \mathbf{j}
        + (u_1v_2 - u_2v_1) \mathbf{k}
\end{align*}
If you are unfamiliar with the operation above, I highly recommend
reading my notes on determinants, which is in Section \ref{sec:26033}!

\begin{important}
The cross product $\mathbf{u} \times \mathbf{v}$ is only defined for
$\mathbf{u}, \mathbf{v} \in \mathbb{R}^3$. The other dimension does
not return meaningful properties and definition of the cross product.
\end{important}

Unlike the dot product, the properties of the cross product are
not so nice.

\begin{theorem}{}{Properties of the Cross Product}
For $\mathbf{u}, \mathbf{v}, \mathbf{w} \in \mathbb{R}^3$ and some
scalar $\lambda$, the following equations hold.
\vspace{-0.8\baselineskip}
\begin{enumerate}
\item
$\mathbf{u} \times \mathbf{v} = - (\mathbf{v} \times \mathbf{u})$
\item
$\mathbf{u} \times \mathbf{u} = \mathbf{0}$
\item
$\mathbf{u} \times \mathbf{0} = \mathbf{0}$
\item
$\mathbf{u} \times (\mathbf{v} + \mathbf{w})
= (\mathbf{u} \times \mathbf{v}) + (\mathbf{u} \times \mathbf{w})$
\item
$\lambda (\mathbf{u} \times \mathbf{v})
= (\lambda \mathbf{u}) \times \mathbf{v}
= \mathbf{u} \times (\lambda \mathbf{v})$
\end{enumerate}
\end{theorem}

There are a lot to prove, but let's start with the first one.

\begin{proof}
Note that by definition, the following equations hold.
\begin{align*}
    \mathbf{u} \times \mathbf{v}
    &= (u_2v_3 - u_3v_2) \mathbf{i} - (u_1v_3 - u_3v_1) \mathbf{j}
        + (u_1v_2 - u_2v_1) \mathbf{k} \\
    \mathbf{v} \times \mathbf{u}
    &= (v_2u_3 - v_3u_2) \mathbf{i} - (v_1u_3 - v_3u_1) \mathbf{j}
        + (v_1u_2 - v_2u_1) \mathbf{k} \\
    &= -(v_3u_2 - v_2u_3) \mathbf{i} + (v_3u_1 - v_1u_3) \mathbf{j}
        - (v_2u_1 - v_1u_2) \mathbf{k}
\end{align*}
Thus, $\mathbf{u} \times \mathbf{v} =
- (\mathbf{v} \times \mathbf{u})$.
\end{proof}

Below is the proof for the second one.

\begin{proof}
By the first property, $\mathbf{u} \times \mathbf{u} =
- (\mathbf{u} \times \mathbf{u})$. In other words,
$2 \mathbf{u} \times \mathbf{u} = \mathbf{0}$ and
$\mathbf{u} \times \mathbf{u} = \mathbf{0}$.
\end{proof}

Continuing, here is the proof for the third property.

\begin{proof}
By definition, $\mathbf{u} \times \mathbf{0} = (u_2v_3 - u_3v_2)
\mathbf{i} - (u_1v_3 - u_3v_1) \mathbf{j} + (u_1v_2 - u_2v_1)
\mathbf{k} = 0\mathbf{i} + 0\mathbf{j} + 0\mathbf{k} = \mathbf{0}$.
\end{proof}

Below is the proof for the fourth property.

\begin{proof}
First and foremost, consider the following equation for the
first component only.
\begin{align*}
    \langle 1, 0, 0 \rangle \cdot
        (\mathbf{u} \times (\mathbf{v} + \mathbf{w}))
    &= u_2(v_3 + w_3) - u_3(v_2 + w_2) \\
    &= (u_2v_3 + u_2w_3) - (u_3v_2 + u_3w_2) \\
    &= u_2v_3 - u_3v_2 + u_2w_3 - u_3w_2 \\
    &= \langle 1, 0, 0 \rangle \cdot \mathbf{u} \times \mathbf{v}
        + \langle 1, 0, 0 \rangle \cdot \mathbf{u} \times \mathbf{w}
\end{align*}
Similarly, the second and the third entry hold with analogous
argument. Thus, the property holds.
\end{proof}

Lastly, here is the proof for the last property.

\begin{proof}
By definition, the following equations hold.
\begin{align*}
    k (\mathbf{u} \times \mathbf{v})
    &= k(u_2v_3 - u_3v_2) \mathbf{i} - k(u_1v_3 - u_3v_1) \mathbf{j}
        + k(u_1v_2 - u_2v_1) \mathbf{k} \\
    &= ((ku_2)v_3 - (ku_3)v_2) \mathbf{i}
        - ((ku_1)v_3 - (ku_3)v_1) \mathbf{j} \\
        &\qquad\qquad\qquad + ((ku_1)v_2 - (ku_2)v_1) \mathbf{k} \\
    &= (k\mathbf{u}) \mathbf{v}
\end{align*}
The analogous computations apply for $k (\mathbf{u} \times \mathbf{v})
= \mathbf{u} \times (k\mathbf{v})$.
\end{proof}

Continuing, it is worth memorizing the cross product between the
standard basis vectors.
\[
\mathbf{i} \times \mathbf{j} = \mathbf{k}, \quad
\mathbf{j} \times \mathbf{k} = \mathbf{i}, \quad
\mathbf{k} \times \mathbf{i} = \mathbf{j}
\]
I personally think and easy way to memorize this is through cycle.
For cycle counter clockwise, we get a negative one.
\[
\mathbf{j} \times \mathbf{i} = -\mathbf{k}, \quad
\mathbf{k} \times \mathbf{j} = -\mathbf{i}, \quad
\mathbf{i} \times \mathbf{k} = -\mathbf{j}
\]
We can also notice that the cross product returns a vector
that are orthogonal to those being crossed!

\begin{important}
Parenthesis is very important! As you can see from the standard
basis vectors above, cross products are generally not associative.
Meaning,
\[
\mathbf{u} \times (\mathbf{v} \times \mathbf{w})
\neq (\mathbf{u} \times \mathbf{v}) \times \mathbf{w}
\]
is true in general and parenthesis is necessarily to avoid ambiguity.
\end{important}

However when we have a dot product and a cross product like
$\mathbf{u} \cdot (\mathbf{v} \times \mathbf{w})$, then parenthesis
is not necessary as the cross product is defined only between two
three dimensional vectors. Speaking of such forms, consider the
following theorem.

\begin{theorem}{}{Cross Product Geometry 1}
For two nonzero vectors $\mathbf{u}$ and $\mathbf{v}$ in the third
dimension, $\mathbf{u} \cdot \mathbf{u} \times \mathbf{v}$ and
$\mathbf{v} \cdot \mathbf{u} \times \mathbf{v}$ holds.
\end{theorem}
\begin{proof}
First, let $\mathbf{u} = \langle u_1, u_2, u_3 \rangle$ and
$\mathbf{v} = \langle v_1, v_2, v_3 \rangle$ and consider the
following equations.
\begin{align*}
    \mathbf{u} \cdot \mathbf{u} \times \mathbf{v} 
    &= \mathbf{u} \cdot \langle
        u_2v_3 - u_3v_2, u_3v_1 - u_1v_3, u_1v_2 - u_2v_1
    \rangle \\
    &= u_1 (u_2v_3 - u_3v_2) + u_2 (u_3v_1 - u_1v_3)
        + u_3 (u_1v_2 - u_2v_1) \\
    &= u_1u_2v_3 - u_1u_3v_2 + u_2u_3v_1 - u_1u_2v_3
        + u_1u_3v_2 - u_2u_3v_1 \\
    &= 0
\end{align*}
Therefore, the first equation holds. Continuing, the following
equations hold for the second equation by Theorem \ref{theo:Properties
of the Cross Product}.
\[
\mathbf{v} \cdot \mathbf{u} \times \mathbf{v}
= \mathbf{v} \cdot (-\mathbf{v} \times \mathbf{u})
= - \mathbf{v} \cdot \mathbf{v} \times \mathbf{u}
= 0
\]
Thus the theorem holds.
\end{proof}

From the theorem, we can notice that $\mathbf{u}$ is orthogonal
to $\mathbf{u} \times \mathbf{v}$, which is orthogonal to
$\mathbf{v}$. These types of operation is known as scalar
triple product, and we will be discussing more on that soon!
In addition to the orthogonal nature of $\mathbf{u} \times \mathbf{v}$
and its adherence to the right-hand rule, we can also find the
relationship between the norm of $\mathbf{u} \times \mathbf{v}$
and $\mathbf{u}$ and $\mathbf{v}$.

\begin{theorem}{}{Cross Product Geometry 2}
For two nonzero vectors $\mathbf{u}$ and $\mathbf{v}$ in the third
dimension with angle $\theta$, $\| \mathbf{u} \times \mathbf{v} \|
= \| \mathbf{u} \| \| \mathbf{v} \| \sin\theta$.
\end{theorem}
\begin{proof}
Let $\mathbf{u} = \langle u_1, u_2, u_3 \rangle$ and
$\mathbf{v} = \langle v_1, v_2, v_3 \rangle$. Notice that because
both the left-hand side and the right-hand side of the equation
are both positive, it suffices to show that the square of each
side equals. Consider the following equations.
\begin{align*}
    \| \mathbf{u} \times \mathbf{v} \|
    &= (u_2v_3 - u_3v_2)^2 + (u_3v_1 - u_1v_3)^2
        + (u_1v_2 - u_2v_1)^2 \\
    &= u_2^2v_3^2 + u_3^2v_2^2 - 2u_2v_3u_3v_2
        + u_3^2v_1^2 + u_1^2v_3^2 - 2u_3v_1u_1v_3 \\
    &\qquad\quad+ u_1^2v_2^2 + u_2^2v_1^2 - 2u_1v_2u_2v_1 \\
    &= u_1^2v_2^2 + u_2^2v_3^2 + u_3^2v_1^2
        + v_1^2u_2^2 + v_2^2u_3^2 + v_3^2u_1^2 \\
    &\qquad\quad- 2u_1u_2v_1v_2 - 2u_2u_3v_2v_3 - 2u_3u_1v_3v_1
\end{align*}
Continuing, the right-hand side can be rearranged as the following.
\begin{align*}
    &\quad\ \| \mathbf{u} \|^2 \| \mathbf{v} \|^2 \sin^2\theta \\
    &= \| \mathbf{u} \|^2 \| \mathbf{v} \|^2
        - \| \mathbf{u} \|^2 \| \mathbf{v} \|^2 \cos^2\theta \\
    &= \| \mathbf{u} \|^2 \| \mathbf{v} \|^2
        - (\mathbf{u} \times \mathbf{v})^2 \\
    &= (u_1^2 + u_2^2 + u_3^2)(v_1^2 + v_2^2 + v_3^2)
        - (u_1v_1 + u_2v_2 + u_3v_3)^2 \\
    &= u_1^2v_1^2 + u_1^2v_2^2 + u_1^2v_3^2
        + u_2^2v_1^2 + u_2^2v_2^2 + u_2^2v_3^2
        + u_3^2v_1^2 + u_3^2v_2^2 + u_3^2v_3^2 \\
    &\qquad - (u_1^2v_1^2 + u_2^2v_2^2 + u_3^2v_3^2)
        - 2(u_1u_2v_1v_2 + u_2u_3v_2v_3 + u_3u_1v_3v_1) \\
    &= u_1^2v_2^2 + u_2^2v_3^2 + u_3^2v_1^2
        + v_1^2u_2^2 + v_2^2u_3^2 + v_3^2u_1^2 \\
    &\qquad\quad- 2u_1u_2v_1v_2 - 2u_2u_3v_2v_3 - 2u_3u_1v_3v_1
\end{align*}
Thus, the theorem holds.
\end{proof}

This theorem implies that the norm of the cross product of the
two vectors is the product of the norms of the vectors and sine of
the angle between them. With this, we can also find another geometric
representation of cross products.

\begin{important}
The area of the parallelogram with sides $\mathbf{u}$ and $\mathbf{v}$
can be represented as $\| \mathbf{u} \times \mathbf{v} \|$.
\end{important}

This is somewhat simple to notice as if we consider $\mathbf{u}$
as the base, then $\| \mathbf{v} \| \sin\theta$ is the height
where $\theta$ is the angle between them. Therefore by Theorem
\ref{theo:Cross Product Geometry 2}, the area of the parallelogram
is $\| \mathbf{u} \| \| \mathbf{v} \| \sin\theta =
\| \mathbf{u} \times \mathbf{v} \|$.

Before we move on to the next section, let's wrap up our discussion
on vectors with the scalar triple product and its geometric
representations.

\begin{definition}{}{}
The \textbf{scalar triple product} of three vectors $\mathbf{u}$,
$\mathbf{v}$, and $\mathbf{w}$ in 3-space is defined as
$\mathbf{u} \cdot \mathbf{v} \times \mathbf{w}$.
\end{definition}

Notice that the output of the scalar triple product of three
vectors is a scalar. Also, using the definition of cross products
and dot products, we can notice that the scalar triple product
is equivalent to the following.
\[
\mathbf{u} \cdot \mathbf{v} \times \mathbf{w}
=
\begin{bmatrix}
    u_1 & u_2 & u_3 \\
    v_1 & v_2 & v_3 \\
    w_1 & w_2 & w_3
\end{bmatrix}
\]
From the matrix operation above, we can notice that the following
equations hold.
\begin{align*}
    \mathbf{u} \cdot \mathbf{v} \times \mathbf{w}
    &= \begin{bmatrix}
        u_1 & u_2 & u_3 \\
        v_1 & v_2 & v_3 \\
        w_1 & w_2 & w_3
    \end{bmatrix}
    = -\begin{bmatrix}
        u_1 & u_2 & u_3 \\
        w_1 & w_2 & w_3 \\
        v_1 & v_2 & v_3
    \end{bmatrix}
    = \begin{bmatrix}
        w_1 & w_2 & w_3 \\
        u_1 & u_2 & u_3 \\
        v_1 & v_2 & v_3
    \end{bmatrix} \\
    &= \mathbf{w} \cdot \mathbf{u} \times \mathbf{v}
\end{align*}
Therefore, we could see that $\mathbf{u} \cdot \mathbf{v} \times
\mathbf{w} = \mathbf{w} \cdot \mathbf{u} \times \mathbf{v}$ and
$\mathbf{w} \cdot \mathbf{u} \times \mathbf{v} = \mathbf{v} \cdot
\mathbf{w} \times \mathbf{u}$ both holds by cycling through
the vectors.

\begin{important}
For any vectors $\mathbf{u}, \mathbf{v}, \mathbf{w}$ in 3-space,
the following equation holds.
\[
\mathbf{u} \cdot \mathbf{v} \times \mathbf{w}
= \mathbf{v} \cdot \mathbf{w} \times \mathbf{u}
= \mathbf{w} \cdot \mathbf{u} \times \mathbf{v}
\]
Moreover because $\mathbf{u} \cdot \mathbf{v} \times \mathbf{w}
= \mathbf{w} \cdot \mathbf{u} \times \mathbf{v}
= \mathbf{u} \times \mathbf{v} \cdot \mathbf{w}$, the following
equation also holds.
\[
\mathbf{u} \cdot \mathbf{v} \times \mathbf{w}
= \mathbf{u} \times \mathbf{v} \cdot \mathbf{w}
\]
\end{important}

Lastly before wrapping up, let's take a look at geometric
interpretation of the scalar triple product. First, consider
a parallelepiped with edges $\mathbf{u}, \mathbf{v}, \mathbf{w}$.

From our previous observation, we can see that the area of the
base of the parallelepiped can be represented as $\| \mathbf{v}
\times \mathbf{w} \|$. Moreover the height can be represented
as $\| \mathbf{u} \| | \cos\theta |$ where $\theta$ is the
angle between $\mathbf{u}$ and $\mathbf{v} \times \mathbf{w}$.
Therefore, the volume can be represented as $\| \mathbf{u} \|
| \cos\theta | \| \mathbf{v} \times \mathbf{w} \| = | \mathbf{u}
\times \mathbf{v} \times \mathbf{w} |$. Continuing with this
idea, we can establish an important theorem.

\begin{theorem}{}{Scalar Triple Product and Coplanar}
$\mathbf{u}, \mathbf{v}, \mathbf{w}$ are coplanar if and only
if $\mathbf{u} \cdot \mathbf{v} \times \mathbf{w} = 0$.
\end{theorem}
\begin{proof}
From the previous observation on the relationship between
the scalar triple product and the volume of the parallelepiped
with the edges $\mathbf{u}, \mathbf{v}, \mathbf{w}$, it is
evident that the volume equals zero if and only if
$\mathbf{u} \cdot \mathbf{v} \times \mathbf{w} = 0$. Moreover,
the volume equals zero if and only if the vectors are coplanar.
\end{proof}

This is it for the vector section! There were lots of topics and
geometric interpretations of vectors that were not included in the
notes for linear algebra. In our next section, we will continue
our discussion through more extensions of geometric representations
of vectors.

\end{document}


